本文围绕K-Waay原有的批内参与方互异条件，构建共享两槽生命周期的无互异约束、消息互异约束和参与方互异约束三种批次准入模型，并以Tamarin进行受控比较。无互异约束模型中，同一匹配发送来源可以对应同批两个接受事件；消息互异约束模型虽同时保持批内消息互异性和批内来源单射性，同一参与方的不同发送实例及不同消息组成的批次仍然可达；作为目标对照，参与方互异约束模型保持批内参与方互异性，且有效的不同参与方批次仍可到达接受阶段。

上述结果表明，消息或发送实例层面的约束不能替代参与方层面的批内关系。若实现需要满足K-Waay既有的批内参与方互异条件，批次构造或准入过程需要依据适当的参与方表示维护相应关系；具体维护机制及其实现映射仍需结合完整协议与系统边界确定。
